% ============================================================
% code-style.tex
% A5 技術書テンプレート用：listings + tcolorbox のコード環境
% ============================================================
% listing の枠線・背景を消す
%	\floatstyle{plain}
%	\restylefloat{listing}

% listing の上下余白を詰める
%	\setlength{\intextsep}{0pt}
%	\setlength{\textfloatsep}{0pt}

% ------------------------------------------------------------
% 必要パッケージ
% ------------------------------------------------------------
% \usepackage{listingsutf8}   % 日本語対応 listings
% \usepackage[most]{tcolorbox}
% \usepackage{xcolor}
\lstset{inputencoding=utf8} 

% ------------------------------------------------------------
% 色定義（コード用）
% ------------------------------------------------------------
\definecolor{codebg}{HTML}{F5F5F5}
\definecolor{keywordblue}{HTML}{005CC5}
\definecolor{stringred}{HTML}{D73A49}
\definecolor{commentgray}{HTML}{777777} % ← 行番号・コメントの色
\definecolor{framegray}{HTML}{D0D7DE}

% ------------------------------------------------------------
% listings 共通スタイル
% ------------------------------------------------------------
\lstdefinestyle{mypython-base}{
  language=Python,
  backgroundcolor=\color{codebg},
  % basicstyle=\codejfont\ttfamily\small,   % ← 和文＋欧文ゴシック体
  basicstyle=\ttfamily\small,
  keywordstyle=\color{keywordblue},
  stringstyle=\color{stringred},
  commentstyle=\codejfont\color{commentgray},
  showstringspaces=false,
  breaklines=true,
  tabsize=4,
  columns=fullflexible,
}

% ------------------------------------------------------------
% 行番号あり
% ------------------------------------------------------------
\lstdefinestyle{mypython}{
  style=mypython-base,
  numbers=left,
  numberstyle=\codejfont\color{commentgray}\small, % ← 太字を外して控えめに
  numbersep=5pt,   % 行番号とコードの距離
}

% ------------------------------------------------------------
% 行番号なし
% ------------------------------------------------------------
\lstdefinestyle{mypython-nonum}{
  style=mypython-base,
  numbers=none,
}

% ------------------------------------------------------------
% tcolorbox によるコードブロック定義
% ------------------------------------------------------------

% 行番号あり：開始行をオプションで指定（デフォルト 1）
\newtcblisting{pycode}[1][1]{%
  listing only,
  listing style=mypython,
  colback=codebg,
  colframe=framegray,
  boxrule=0.4pt,
  arc=3pt,
  outer arc=3pt,
  left=2.0em,   % ← 行番号あり／なしで左端を揃える
  right=4pt,
  top=2pt,
  bottom=2pt,
  enhanced,
  breakable,
  before upper=\lstset{firstnumber=#1},
}

% 行番号なし
\newtcblisting{pycodenonum}{%
  listing only,
  listing style=mypython-nonum,
  colback=codebg,
  colframe=framegray,
  boxrule=0.4pt,
  arc=3pt,
  outer arc=3pt,
  left=2.0em,   % ← 行番号なしでも左端を揃える
  right=4pt,
  top=4pt,
  bottom=4pt,
  enhanced,
  breakable,
}

% ============================================================
% ここまで
% ============================================================
